\Needspace{9\baselineskip}
\section{Nonidentities and the uniform height bound}
\label{nid:section}

We verify the last hypothesis of \Cref{int:height}: the equations selected
in \Cref{loc:groups} cannot all vanish identically on any algebraic branch
that occurs in the interpolation argument.  This assertion concerns
arbitrary algebraic subvarieties of the coordinate space.  Their
endomorphism coordinates are not required to represent Hodge classes.
Monodromy treats the single and paired equations.  For the exceptional
equation it leaves two signs.  A Frobenius multiplier excludes the
positive sign; the degenerate first auxiliary incidence excludes the
negative sign through the integral Tate lattice.

\subsection{Algebraic branches and continuation}

Fix one group from \Cref{loc:groups}, its target $z=(a,U)$ with $a\ne0$,
and its equations, including the base-coordinate equations.  Center data,
conjugate families, and any matrix $r$ or isometry $I_0$ are constants in
these equations.  The operator coordinates and, when used, $R$ are
coordinates of the ambient affine space.  Write $\mathscr E$ for this
finite collection of analytic expressions.  Its algebraic function germs
are taken on the specified branches through the centers.

\begin{lemma}[Continuation from a normalized branch]
\label{nid:continuation}
Let $W$ be an irreducible closed algebraic subvariety of the ambient
affine space, containing $z$, on which $a$ is nonconstant.  Let $D$ be a
prime divisor on the normalization of $W$ lying over the generic point
of a component of $W\cap\{a=0\}$.  If every expression in $\mathscr E$
vanishes as a germ on this branch, then the following hold.
\begin{enumerate}[label=\textup{(\roman*)}]
\item The base coordinates give actual algebraic projections to the
fixed curve charts; the first projection $s$ is nonconstant and satisfies
$x(s)=a$.  In the paired case the second satisfies $x'(s')=aR$.
\item On a connected smooth algebraic open $W^\circ$ all the identities
continue along every path.  At a fixed point of $W^\circ$, continuation
around a loop changes the horizontal transports by their pullback
monodromy, while returning the algebraic coordinates to their original
values.
\item Every operator pair that is independent at $z$ is independent at
a general point of $W^\circ$.
\end{enumerate}
The same conclusions hold if the pullbacks of all expressions vanish to
infinite order along $D$.
\end{lemma}

\begin{proof}
At a general complex point of $D$ the normalization and $D$ are smooth.
Take local coordinates $(u,v_2,\ldots,v_k)$ there with $D=\{u=0\}$.
The regular function $a$ is $u^e$ times a unit, for some $e>0$.
All affine coordinates of $W$ pull back to holomorphic functions; in
particular $R$ is finite there.  Consequently both $a$ and $aR$ tend to
zero, and the ordinary-point series defining our expressions converge
on a sufficiently small neighborhood.  An infinite-order zero along
$D$ makes every coefficient of their convergent expansion in $u$ zero.
Thus it gives an identity on the whole local branch.  This is stronger
than an identity obtained by setting $u=0$.

Let $c_\alpha(x)$ be the algebraic coordinate germs of the first curve
chart.  Its added equations are $s_\alpha-c_\alpha(a)=0$.  Every
polynomial relation defining the chart, and the relation $x(s)=a$,
therefore holds on the analytic branch.  Such a relation is an algebraic
identity on $W$, since a nonempty analytic open of an irreducible complex
algebraic variety is Zariski dense.  This constructs the asserted
projection to the actual chart; the same argument applies to
$s'_\alpha-c'_\alpha(aR)$ when present.  Relations involving regular
functions on these affine charts can equivalently be checked after
clearing denominators that are invertible on the chart.  The target
belongs to these charts by \Cref{loc:equations}, so the algebraic
relations also hold at the specified target.  Nonconstancy of $a$
implies nonconstancy of $s$.

Remove singularities and the proper closed sets where the chosen frames
or projections are undefined.  Shrink further to avoid the branch loci
of any parameter charts used to describe the germs.  A nonempty part of
our analytic branch lies over the resulting smooth open: normalization
is an isomorphism over the smooth locus of $W$.  This open is connected.
Solutions of the pulled-back connections continue along its paths, and
the identity theorem preserves the equations under this continuation.
All their algebraic coordinate functions are single-valued there.  This
proves (ii), without requiring the transports themselves to be
single-valued.

Finally, an independent pair at $z$ has a two-by-two coordinate minor
nonzero at $z$.  That minor is a nonzero regular function on $W$.
Deleting its zero set for each of the finitely many pairs proves (iii).
No Hodge condition on these varying coordinates has entered the proof.
\end{proof}

For use below, a nonconstant projection from $W^\circ$ to a fixed base
curve has the full connected monodromy of that curve.  Indeed, choose
an algebraic curve in $W^\circ$ on which the projection is nonconstant.
After normalization and omission of finitely many points it is a finite
cover of an open of the base curve.  Its monodromy has finite index in
the monodromy on that open, and hence the same connected Zariski closure.
The monodromy of $W^\circ$ contains this image.  The joint assertion for
two nonconstant projections is precisely \Cref{geo:monodromy}; its
exceptional images are contained in the fixed closed correspondences
of \Cref{geo:hecke-list}.

\begin{lemma}[Matrix coefficients and their span]
\label{nid:matrix-span}
Let $V$ be the complex five-dimensional quadratic representation of
\Cref{geo:five-space}, and let $G=\SO(V)$.
\begin{enumerate}[label=\textup{(\roman*)}]
\item If $b,v\in V$ are nonzero, the function
$g\mapsto\langle b,gv\rangle$ on $G$ is not constant.
\item The complex linear span of $G$ is $\End_{\C}(V)$.
Consequently, if $b,c\in V$, $p,q\in V_1$, and $M_0:V_1\to V$ is
an isomorphism, the identities
\[
 \langle b,gM_0p\rangle=\langle c,gM_0q\rangle\qquad(g\in G)
\]
imply $b\otimes p=c\otimes q$, where the first tensor factors are
identified with their duals by the quadratic pairing.
\end{enumerate}
The assertions remain valid if the identities are initially known only
on a Zariski-dense monodromy subgroup.
\end{lemma}

\begin{proof}
For (i), let $V'$ be the span of $(g-1)v$ for $g\in G$.  It is invariant,
since
\[
 h(g-1)v=(hg-1)v-(h-1)v.
\]
It is nonzero, as the nontrivial irreducible representation $V$ has no
fixed nonzero vector.  Irreducibility gives $V'=V$.  Constancy of the
displayed coefficient would make $b$ perpendicular to $V'$, a
contradiction.

Assertion (ii) is the matrix-span conclusion of Burnside's theorem; in
this representation it also has the following direct verification.
Choose an orthonormal basis.  The diagonal matrices
$D_\eta=\diag(\eta_1,\ldots,\eta_5)$ with $\eta_i\in\{1,-1\}$ and
$\prod_i\eta_i=1$ lie in $G$.  Their coordinate characters are distinct,
so
\[
 E_{ii}=\frac1{16}\sum_{\prod_k\eta_k=1}\eta_iD_\eta
\]
lies in their span.  The span of a matrix group is an algebra.  Signed
permutation matrices of determinant one, together with the $E_{ii}$,
therefore give every matrix unit $E_{ij}$ in this algebra.  This proves
(ii).  The displayed identity is a linear functional of $gM_0$;
nondegeneracy of the tensor--matrix pairing proves the tensor conclusion.
Polynomial identities on a Zariski-dense subgroup extend to $G$.
\end{proof}

\subsection{The Frobenius obstruction}

For the exceptional equation, a hypothetical analytic identity will
give four tensor equalities by \Cref{nid:matrix-span}.  At the target,
these force one common scalar $\lambda$: the center pair is multiplied
by $\lambda$, and the target pair by $\lambda^{-1}$.  The center
endomorphisms have nonzero rational trace norms, unchanged by
conjugation and by the isometry $I_0$, so $\lambda^2=1$.
We isolate the arithmetic obstruction to these two signs first.
The lemma concerns actual endomorphisms of the characteristic-zero
fibers; we will apply it only after deriving their equalities at the
target from the analytic identity.

\begin{lemma}[Exclusion of both signs]
\label{nid:frobenius}
In an exceptional group of \Cref{loc:equations}, the target data cannot
satisfy, for either common sign $\epsilon\in\{1,-1\}$,
\begin{equation}
 \label{nid:sign-equalities}
 I_0^{-1}B'_i=\epsilon B_i\quad(i=1,2),
 \qquad
 I_t^{-1}P'_j=\epsilon P_j\quad(j=1,2)
\end{equation}
as equalities of rational endomorphisms.  Here $I_0$ and $I_t$ are
conjugation by the fixed-multiplier polarized isogenies specified in
\Cref{loc:equations}.
\end{lemma}

\begin{proof}
Choose any place in this group, of residue characteristic $p$.
\Cref{loc:equations} gives a common unprimed polarized reduction $A_0$
and its primed reduction $A'_0$, a relative Frobenius
$F:A_0\to A'_0$ of multiplier $p$, and a polarized isogeny
$\phi:A_0\to A'_0$ of multiplier $m_*$ which is simultaneously the
specialization of the center and target isogenies.  In particular,
$p\nmid m_*$ and $p\ne \ell_1$.  Set
\begin{equation}
 \label{nid:psi}
 \psi=\phi^{-1}F\in\End(A_0)\otimes\Q,
 \qquad \nu(\psi)=\frac{p}{m_*},
 \qquad g=\Ad(\psi)\big|_{\cP_{\Q_{\ell_1}}}.
\end{equation}
Functoriality of specialization preserves
\eqref{nid:sign-equalities}.  Moreover $F$ conjugates each unprimed
labeled endomorphism to its primed specialization, for both the center
and target labels using this same $F$.  Composing with $\phi^{-1}$
therefore shows that $g$ acts by $\epsilon$ on each of the two
specialized rational endomorphism planes.

Let $H_0,H_t$ be these planes in $\End(A_0)\otimes\Q$ and put
$H=H_0+H_t$.  On the common integral Tate module, write
\[
 \begin{gathered}
 \Lambda=\{D\in\End_{\Z_{\ell_1}}(T_{\ell_1}A_0):
                    D^\dagger=D,\ \Tr(D)=0\},\\
 L_0=(H_0\otimes\Q_{\ell_1})\cap\Lambda,
 \qquad L_t=(H_t\otimes\Q_{\ell_1})\cap\Lambda.
 \end{gathered}
\]
The two plane lattices here are individually saturated in $\Lambda$.
Their reductions inject into $\Lambda/\ell_1\Lambda$ and, by
\Cref{inc:markings,loc:equations}, are independent planes whose sum is
a degenerate four-dimensional subspace.  In particular the two
$\Q_{\ell_1}$-planes have four-dimensional sum, and $\dim_\Q H=4$.
These statements use the saturated planes, regardless of whether the
chosen small pairs $\mathbf B,\mathbf P$ are integral bases at $\ell_1$.

Every element of $H$ is an actual rational endomorphism of $A_0$,
self-adjoint for its Rosati involution.  The endomorphism trace is
rational and agrees with its Tate-module trace
\cite[Proposition~12.9]{MilneAV}.  Rosati positivity in arbitrary
characteristic \cite[Theorem~17.3]{MilneAV} gives
\[
 \Tr(D^2)=\Tr(DD^\dagger)>0\qquad(0\ne D\in H).
\]
Thus the Gram matrix of a rational basis of $H$ is positive definite
over $\R$, with nonzero rational determinant.  It remains nonsingular
over $\Q_{\ell_1}$.  Write $H_{\ell_1}=H\otimes\Q_{\ell_1}$.  This rational
nondegeneracy is compatible with the prescribed degeneracy of the
reduced four-space: its Gram determinant may be divisible by $\ell_1$.

Conjugation by a symplectic similitude acts through $\SO$ on $\cP$,
with scalar kernel on the four-dimensional first homology space, by
\Cref{geo:five-space}.  Since $H_{\ell_1}$ is nondegenerate, $g$ preserves
its one-dimensional orthogonal complement.  On $H_{\ell_1}$ its
determinant is $\epsilon^4=1$; hence $\det g=1$ makes its action on the
orthogonal complement $+1$.

If $\epsilon=1$, this proves $g=1$ on all of $\cP_{\Q_{\ell_1}}$.
The scalar-kernel assertion makes the action of $\psi$ on
$T_{\ell_1}A_0\otimes\Q_{\ell_1}$ a scalar $c$.  This scalar is rational:
the actual rational quasi-endomorphism $\psi$ has rational trace and
$4c=\Tr(\psi)$, by \cite[Proposition~12.9]{MilneAV}.  Its symplectic
multiplier is consequently $c^2$.  Equation~\eqref{nid:psi} would give
$c^2=p/m_*$, impossible because the $p$-adic valuation of the right
side is one.

Suppose instead that $\epsilon=-1$.  Then $g^2=1$ and its minus
eigenspace is exactly $H_{\ell_1}$.  Both $F$ and $\phi$ have degrees prime
to $\ell_1$: their degrees are $p^2$ and $m_*^2$, respectively.
They induce integral isomorphisms on the $\ell_1$-adic Tate modules.
Consequently $\psi$, its inverse, and their conjugation action preserve
$\Lambda$.  The trace pairing on $\Lambda$ is unimodular, as in
\Cref{geo:five-space} and the choice of $\ell_1$ in
\Cref{inc:markings}.  Since $\ell_1$ is odd, the integral idempotents
\[
 e_+=\frac{1+g}{2},\qquad e_-=\frac{1-g}{2}
\]
give an orthogonal direct sum
\begin{equation}
 \label{nid:lattice-splitting}
 \Lambda=\Lambda_+\mathbin{\perp}\Lambda_-,
 \qquad \Lambda_\pm=e_\pm\Lambda,
 \qquad \rank\Lambda_-=4.
\end{equation}
Orthogonality follows from
$\langle v_+,v_-\rangle=\langle gv_+,gv_-\rangle
=-\langle v_+,v_-\rangle$.
In bases of these two direct summands the Gram determinant of $\Lambda$
is the product of their integral Gram determinants.  It is a unit, so
both determinants are units.  In particular
$\Lambda_-/\ell_1\Lambda_-$ is nondegenerate.

Because $g=-1$ on each rational plane, $L_0$ and $L_t$ lie in
$\Lambda_-$.  Their independent reductions have total dimension four.
The direct-summand property in \eqref{nid:lattice-splitting} identifies
$\Lambda_-/\ell_1\Lambda_-$ with its image in $\Lambda/\ell_1\Lambda$.
Those two reductions therefore exhaust this image.  Their sum must be
nondegenerate, contrary to its imposed degenerate incidence.  This
excludes the second sign as well.
\end{proof}

\subsection{The three equation types}

\begin{proposition}[Nonidentity on every algebraic branch]
\label{nid:nonidentity}
For every group selected in \Cref{loc:groups}, every irreducible closed
algebraic $W$ containing its target and having nonconstant $a$, and every
normalized branch over the generic point of a component of
$W\cap\{a=0\}$, at least one of its equations is a nonzero germ on that
branch.  This holds for each of the three types
\eqref{loc:single}, \eqref{loc:paired}, and \eqref{loc:exception}, with
their base-coordinate equations included.
\end{proposition}

\begin{proof}
Suppose all equations were identities on one such branch.  Apply
\Cref{nid:continuation}, and fix a general point $w\in W^\circ$ at which
the indicated operator pairs are independent.  All vectors and
coordinates at $w$ are now fixed while we vary a continuation loop.
The continued first transport $M:\cP_{s(w)}\to\cP_{t_0}$ ranges through
the orbit of an invertible initial transport under monodromy whose
connected Zariski closure is the full special orthogonal group.

\paragraph{The single equation.}
Equation~\eqref{loc:single} says, for all $i,j\in\{1,2\}$,
\[
 \langle B_i,MP_j(w)\rangle=r_{ij}.
\]
The right side is fixed under continuation.  Each $B_i$ and $P_j(w)$
is nonzero, and the first connected monodromy is full.  This contradicts
\Cref{nid:matrix-span}(i).

\paragraph{The paired equation.}
Here the equation is
\[
 \langle B_i,MP_j(w)\rangle
   =\langle B'_i,M'P'_j(w)\rangle.
\]
If the second curve projection is constant, its pullback local system
has trivial monodromy, so varying the first connected monodromy already
contradicts \Cref{nid:matrix-span}(i).  If both projections are
nonconstant, \Cref{geo:monodromy,geo:hecke-list} give full product
monodromy unless the joint image lies in one of the fixed closed Hecke
correspondences.  In that exception its algebraic closure also contains
the target pair in that correspondence: the coordinate projections are
defined at the target and the correspondence is closed.  This is
excluded by the definition of a paired group in \Cref{loc:equations}.
We therefore have product monodromy.  The equations are polynomial in
the two transport matrices and hence hold on its Zariski closure.
Hold the second factor fixed and vary the first; the same nonconstant
matrix coefficient gives a contradiction.

\paragraph{The exceptional equation.}
Put $C_i=I_0^{-1}B'_i$, which is a constant center vector.  The coordinate
pair $\widehat{\mathbf P}$ equals $I_t^{-1}\mathbf P'$ at the target;
away from that target it is simply a pair of algebraic operator
coordinates.  Equation~\eqref{loc:exception} and full first monodromy
give
\[
 \langle B_i,MP_j(w)\rangle
   =\langle C_i,M\widehat P_j(w)\rangle
       \qquad(i,j\in\{1,2\})
\]
for the entire monodromy orbit of $M$.  By
\Cref{nid:matrix-span}(ii),
\begin{equation}
 \label{nid:tensors}
 B_i\otimes P_j=C_i\otimes\widehat P_j
       \qquad(i,j\in\{1,2\})
\end{equation}
at a general point of $W$.  In the algebraic frames these are algebraic
identities in the operator coordinates.  Thus they hold everywhere on
$W$, in particular at $z$.

At the target all vectors in the displayed pairs are nonzero.
The equality for $(i,j)=(1,1)$ gives a scalar $\lambda\ne0$ with
$C_1=\lambda B_1$ and
$\widehat P_1=\lambda^{-1}P_1$.  The equality for $(2,1)$ then gives
$C_2=\lambda B_2$, and that for $(1,2)$ gives
$\widehat P_2=\lambda^{-1}P_2$.  This is one common scalar for all four
cross entries.

Conjugation by $I_0$ preserves the trace pairing.  Moreover each fixed
center norm $\langle B_i,B_i\rangle$ is a positive rational number,
since $B_i$ is an actual nonzero Rosati-self-adjoint rational
endomorphism.  Conjugation of the number field preserves this rational
trace, so
\[
 \langle C_i,C_i\rangle
  =\langle B'_i,B'_i\rangle
  =\langle B_i,B_i\rangle\ne0.
\]
It follows that $\lambda^2=1$.  Write $\epsilon=\lambda$; then
$\lambda^{-1}=\epsilon$ also, and the target equalities are exactly
\eqref{nid:sign-equalities} in de Rham realization.

They are equalities of actual rational endomorphisms.  Indeed all the
target vectors used here represent such endomorphisms, and de Rham
realization in characteristic zero is faithful: after embedding the
field in $\C$, comparison gives their action on first Betti homology;
a homomorphism of complex abelian varieties acting trivially on that
lattice is zero.  Clearing rational denominators gives the same
assertion for quasi-endomorphisms.  We may therefore apply
\Cref{nid:frobenius}, which excludes both values of $\epsilon$.
This finishes all three cases.
\end{proof}

\Needspace{12\baselineskip}
\subsection{Uniform height conclusion}

\begin{theorem}[Height of the selected lifts]
\label{nid:height}
There are constants $C>0$ and $\kappa>0$, depending only on the fixed
curve, covers, center, and auxiliary data, such that every selected
non-CM elliptic-square lift $t$ satisfies
\begin{equation}
 \label{nid:height-bound}
 h(t)\le C\max\{2,[K(t):K]\}^{\kappa}.
\end{equation}
The height is the shifted ample height of \Cref{ht:conventions}.
\end{theorem}

\begin{proof}
Put $d=\max\{2,[K(t):K]\}$ and $h=h(t)\ge2$.  All fixed omitted
points may be absorbed by increasing $C$.  For the remaining targets,
\Cref{loc:groups} either already gives a polynomial bound for $h$ in
$d$, or supplies a target $z=(a,U)$, a weighted group $V$, and common
equations to which we apply \Cref{int:height}.

We match its hypotheses explicitly.  The target has $a\ne0$ and
$|a|_v<1$ on $V$.  The global height, proximity mass, and positive
coordinate bounds in
\eqref{loc:group-height-mass}--\eqref{loc:group-local-degree} give
condition \textup{(i)}.  The coefficient bounds and common tail
estimate in \eqref{loc:coefficient-height}--\eqref{loc:tail-bound},
together with the degree bound in \eqref{loc:group-local-degree},
give \textup{(ii)}.  The expressions are polynomial combinations of
fixed bounded degree of ordinary horizontal systems and regular
algebraic germs, evaluated at $a$ and, in the paired case, at $aR$;
this is \textup{(iii)}.  Here $|R|_v=1$ on the selected group.
The base-coordinate equations are included in this system.

These estimates use the same $d=\max\{2,[K(t):K]\}$ for the fixed
field $K$.  The selected labeled data and isogenies have a field of
definition of degree polynomial in $d$.  The normalized weighted sums,
including embedding multiplicities, are preserved by every further
field extension.  The three equation
types involve only finitely many fixed choices of ambient dimension,
germs, conjugate centers, and center isogenies.  We may therefore choose
the constants in \Cref{int:height} uniformly across them.

Finally, \Cref{nid:nonidentity} gives condition \textup{(iv)} for
every irreducible algebraic variety through $z$ and every normalized
boundary branch, including those arising in the dimension descent.
Varieties with empty boundary are already handled in the proof of
\Cref{int:height}.  That theorem gives $h\le Cd^\kappa$; taking the
largest constants over the fixed alternatives proves
\eqref{nid:height-bound}.
\end{proof}
